\begin{enumerate}[label=\textbf{\arabic*)}, leftmargin=*]

    \item 
    \begin{itemize}
        \item \textbf{transdutor:} Transforma um tipo de energia em outro, podem ser encontrados em estágio de entrada e saída.
        \item \textbf{sensor:} capta uma grandeza física.
        \item \textbf{atuador:} Transforma um sinal elétrico em outra grandeza física.
    \end{itemize}

    \item Um transdutor ativo possui alimentação externa própria, enquanto um transdutor passivo não necessita de energia adicional.

    \item Um sinal analógico possui um sinal contínuo, ou seja, sua curva não possui bicos, já o sinal digital varia entre zeros e uns.

    \item Para a saída em corrente, o instrumento/sensor gera um sinal em forma de corrente elétrica; para saída em tensão, o sensor gera saída em tensão.

    \item Em um instrumento caracterizado como simples há apenas uma conversão no sinal de entrada, no instrumento composto há mais de uma conversão no sinal de entrada.

    \item No método de deflexão a grandeza provoca uma resposta direta (deflexão) no instrumento, e essa resposta é usada para medir a grandeza. No método de detecção de erro nulo, a grandeza é comparada a um sinal de referência ajustável, até que o erro seja nulo.

    \item Os erros aleatórios interferem em cada medida de maneira aleatória, enquanto erros sistemáticos acontecem de maneira semelhante.
    \item \textbf{Análise do alvo (Exatidão / Precisão):}
    \begin{enumerate}[label=\alph*)]
        \item Alta exatidão e alta precisão.
        \item Baixa exatidão e alta precisão.
        \item Alta exatidão e baixa precisão.
        \item Baixa exatidão e baixa precisão.
    \end{enumerate}

    \item \textbf{Algarismos significativos:}
    \begin{enumerate}[label=\alph*)]
        \item $4{,}35 \rightarrow 3$ algarismos significativos.
        \item $0{,}01 \rightarrow 1$ algarismo significativo.
        \item $0{,}0000052 \rightarrow 2$ algarismos significativos.
        \item $400{,}510 \rightarrow 6$ algarismos significativos.
        \item $1000{,}1 \rightarrow 5$ algarismos significativos.
        \item $0{,}00000000000001 \rightarrow 1$ algarismo significativo.
    \end{enumerate}

    \item \textbf{Arredondamento:}
    \begin{enumerate}[label=\alph*)]
        \item $2{,}2\bar{4}5 \rightarrow 2{,}25$
        \item $6\bar{9},87 \rightarrow 70$
        \item $123,\bar{1}2 \rightarrow 123{,}1$
    \end{enumerate}

    \item \textbf{Operações com algarismos significativos:}
    \begin{enumerate}[label=\alph*)]
        \item $12{,}45 \times 7{,}2 = 89{,}64 \approx 90$
        \item $5{,}2 + 7{,}25 = 12{,}45 \approx 12{,}5$
        \item $351{,}12 - 136 = 215{,}12 \approx 215$
        \item $134{,}15 / 2{,}15 = 6{,}2395 \approx 6{,}24$
        \item $(12{,}10 + 5{,}1) \times (6{,}75 - 4{,}89) = 17{,}2 \times 1{,}86 = 31{,}992 \approx 32$
    \end{enumerate}

    \item
    \begin{enumerate}[label=\alph*)]
        \item \textbf{Médias aritméticas:}
        \[
        \bar{P} = \frac{24{,}95 + 25{,}34 + 25{,}14 + 24{,}86 + 25{,}43}{5} = 25{,}14 \approx 25
        \]
        \[
        \bar{i} = \frac{2{,}95 + 3{,}05 + 3{,}00 + 2{,}98 + 3{,}10}{5} = 3{,}02 \approx 3
        \]

        \item \textbf{Desvio padrão da média:}
        \[
        \sigma_P = 0{,}24419 \implies \sigma_{\bar{P}} = \frac{0{,}24419}{\sqrt{5}} = 0{,}1092 = \Delta P
        \]
        \[
        \sigma_I = 0{,}05941 \implies \sigma_{\bar{I}} = \frac{0{,}05941}{\sqrt{5}} = 0{,}02657 = \Delta I
        \]

        \item \textbf{Cálculo da Tensão ($V = P / I$) e Incerteza Propagada:}
        \[
        \bar{V} = \frac{\bar{P}}{\bar{I}} = \frac{25}{3} = 8{,}33\,\text{V} \approx 8\,\text{V}
        \]
        
        Derivadas parciais:
        \[
        \frac{\partial V}{\partial P} = \frac{1}{I} = \frac{1}{3}
        \]
        \[
        \frac{\partial V}{\partial I} = -P \cdot I^{-2} = -25 \cdot 3^{-2} = -2{,}778
        \]

        Propagação da incerteza:
        \[
        \Delta V = \left| \frac{\partial V}{\partial P} \right| \Delta P + \left| \frac{\partial V}{\partial I} \right| \Delta I
        \]
        \[
        \Delta V = \frac{1}{3} \cdot 0{,}11 + 2{,}77 \cdot 0{,}027 = 0{,}112
        \]

        Cálculo do valor médio direto de $V$:
        \[
        \bar{V} = \frac{8{,}46 + 8{,}31 + 8{,}38 + 8{,}34 + 8{,}20}{5} = 8{,}34 \approx 8
        \]
        \[
        \sigma_{\bar{V}} = 0{,}08541
        \]

        \item \textbf{Resultados no formato padrão:}
        \begin{align*}
        V &= 8{,}34 \pm 0{,}09 \\
        i &= 3{,}02 \pm 0{,}05 \\
        P &= 25{,}14 \pm 0{,}01
        \end{align*}
    \end{enumerate}

\end{enumerate}